\documentclass[10pt,a4paper]{article}
\usepackage{amsmath,amssymb,amsthm,graphicx}
\usepackage{xfrac}
% Minimal editorial layout adapter for the pinned JEP source
% sources/jep-2022-195-contact-systolic/pinned-original-native.tex.
% The publisher document class (cedram) is not part of the retained material, and it
% is the class that used to define a handful of plain text abbreviations used by the
% author's prose. This adapter supplies exactly those abbreviations and nothing else:
% no mathematics, no notation, no theorem environment, no counter.
% HoTT/JEP retention convention: CC BY 4.0, attribution to the original authors.
\providecommand{\ie}{i.e.\ }
\providecommand{\eg}{e.g.\ }
\providecommand{\skpt}{\smallskip}
\usepackage[scr=rsfs,cal=euler]{mathalfa}
\let\mathds\mathbb
\let\epsilon\varepsilon
\newcommand{\Changel}{\mathcode`l="7160}
\newcommand{\Changelback}{\mathcode`l="716C}
\newcommand{\NoChangel}[1]{%
\expandafter\let\csname old\string#1\endcsname=#1
\let#1=\relax
\newcommand{#1}{\mathcode`l="716C\csname old\string#1\endcsname\mathcode`l="7160 }%
}

\NoChangel{\log}
\NoChangel{\ln}
\NoChangel{\lim}
\NoChangel{\limsup}
\NoChangel{\liminf}
\NoChangel{\varlimsup}
\NoChangel{\varliminf}
\NoChangel{\varinjlim}
\NoChangel{\varprojlim}

\newenvironment{enumeratei}
{\bgroup\def\theenumi{\roman{enumi}}\def\theenumii{\arabic{enumii}}\begin{enumerate}}
{\end{enumerate}\egroup}
\newcommand{\Psfrac}[2]{(\sfrac{#1}{#2})}
\newcommand{\psfrac}[2]{\sfrac{(#1)}{#2}}
\newcommand{\spfrac}[2]{\sfrac{#1}{(#2)}}
\newcommand\mto{\mathchoice{\longmapsto}{\mapsto}{\mapsto}{\mapsto}}
\newcommand\hto{\mathrel{\lhook\joinrel\to}}

\usepackage{graphicx}

\newtheorem{thm}{Theorem}[section]
\newtheorem{cor}[thm]{Corollary}
\newtheorem{lem}[thm]{Lemma}
\newtheorem{prop}[thm]{Proposition}
\newtheorem{quest}{Question}[section]
\newtheorem{conj}{Conjecture}[section]
\newtheorem{maintheorem}{Theorem}
\theoremstyle{definition}
\newtheorem{rem}[thm]{Remark}
\newtheorem{exm}[thm]{Example}
\newtheorem{ass}[thm]{Assumption}
\renewcommand*{\themaintheorem}{\Alph{maintheorem}}

\DeclareMathOperator{\area}{area}
\DeclareMathOperator{\lcm}{lcm}
\DeclareMathOperator{\const}{const}
\DeclareMathOperator{\Cal}{Cal}
\DeclareMathOperator{\Flux}{\mathrm{Flux}}
\newcommand{\id}{\mathrm{id}}
\newcommand{\Ham}{\mathrm{Ham}}
\newcommand{\std}{\mathrm{std}}
\newcommand{\vol}{\mathrm{vol}}
\newcommand{\fix}{\mathrm{fix}}
\newcommand{\pr}{\mathrm{pr}}
\newcommand{\Tan}{\mathrm{T}}
\newcommand{\N}{\mathds{N}}
\newcommand{\Z}{\mathds{Z}}
\newcommand{\R}{\mathds{R}}
\newcommand{\CC}{\mathds{C}}\let\C\CC
\newcommand{\UU}{\mathcal{U}}
\newcommand{\VV}{\mathcal{V}}

\let\eul e
\let\diff d
\title{The contact volume of a global surface of section is the integral of the first return time}
\author{Alberto Abbondandolo, Christian Lange, Marco Mazzucchelli}
\date{doi:10.5802/jep.195; CC BY 4.0}
\begin{document}
\maketitle
\setcounter{section}{3}
\setcounter{thm}{4}
\noindent\textit{Editorial note.} The counted statement is the author's labelled lem environment at pinned-original-native.tex lines 1436-1442; its proof environment opens at line 1444, so the pairing is by adjacency and the retained proof range 1444-1458 is the whole environment.\par\medskip
\noindent\textit{Editorial note.} Numbering: the lemma carries the author label l:volume and the retained counter values (section 3, theorem offset 4) print it as Lemma 3.5. The number was obtained by counting the theorem-class environments of section 3 against the anchors already fixed for the earlier cards of this source (volume\_Besse 3.2 at line 1227, nearly\_Besse 3.3 at line 1289, return\_time\_map 3.6 at line 1462, pullback\_to\_nu0 3.7 at line 1501), so it is a counter-based inference of the same kind as those, not a value read from a label-number dump.\par\medskip
\noindent\textit{Editorial note.} No retained line refers to a label, citation or figure outside the retained span, so the package declares no external reference.\par\medskip
\noindent\textit{Editorial note.} The retained preamble is the minimal editorial adapter for the publisher's text abbreviations followed by the native author macro and theorem block of the pinned source.\par\medskip
\noindent\textit{Editorial scope.} One original result of the pinned publisher article, counted once: the identity expressing the contact volume as the integral of the first return time over a global surface of section (author label l:volume).\par\medskip
\noindent\textit{Allowed standard prerequisites (editorial).}\begin{itemize}
\item Contact volume. vol(Y,lambda) is the integral over Y of lambda wedge d lambda, for the contact form lambda on the three-manifold Y.. It is the left-hand side of the identity.
\item Global surface of section with first return time. Sigma is a global surface of section for the Reeb flow, tau is the first return time and omega is the area form induced by the contact form on Sigma.. They are the objects whose integral gives the volume.
\item Reeb flow and the map j. The flow phi of the Reeb vector field sends Sigma x [0,1) onto a full measure subset of Y and satisfies j*(lambda wedge d lambda) = tau ds wedge pr\_1\textasciicircum{}* omega.. It is the change of variables that turns the integral over Y into an integral over Sigma.
\end{itemize}\medskip
\begin{lem}
\label{l:volume}
The contact volume of $(Y,\lambda)$ is given by
\[
\vol(Y,\lambda) = \frac1\alpha \int_\Sigma \tau\,\omega.
\]
\end{lem}

\begin{proof}
The bijective map
\[
j: \Sigma \times [0,1) \to \Sigma \times S^1, \qquad j(z,s) := \phi_{\tilde{\lambda}}^{s\tau(z)} (z,0),
\]
satisfies
\[
j^*(\tilde\lambda \wedge d\tilde\lambda) = \tau\, ds \wedge \mathrm{pr}_1^* \omega,
\]
where $\pr_1:\Sigma\times [0,1) \to\Sigma$ is the projection onto the first factor. Together with the fact that the restriction of $\tilde{\iota}$ to the interior of $\Sigma\times S^1$
is an $\alpha$-th fold branched covering map of a full measure subset of $Y$ pulling the volume form $\lambda\wedge d\lambda$ back to $\tilde\lambda \wedge d\tilde\lambda$, we obtain
\begin{align*}
\alpha\, \vol(Y,\lambda) &= \alpha \int_Y \lambda \wedge d\lambda = \int_{\Sigma\times S^1} \tilde\lambda \wedge d\tilde\lambda = \int_{\Sigma\times [0,1)} j^*( \tilde\lambda \wedge d\tilde\lambda) \\ &= \int_{\Sigma\times [0,1)} \tau\, ds \wedge \mathrm{pr}_1^* \omega = \int_{\Sigma} \tau \, \omega.\qedhere
\end{align*}
\end{proof}

\begin{thebibliography}{99}
\bibitem{Abbondandolo:2019tl} A. Abbondandolo \& G. Benedetti - “On the local systolic optimality of Zoll contact forms” , 2019
\bibitem{Abbondandolo:2018fb} A. Abbondandolo, B. Bramham, U. L. Hryniewicz \& P. A. S. Salomão - “Sharp systolic inequalities for Reeb flows on the three-sphere” , Invent. Math. 211 (2018) no. 2, p. 687-778
\bibitem{Abbondandolo:2019ta} A. Abbondandolo, B. Bramham, U. L. Hryniewicz \& P. A. S. Salomão - “Contact forms with large systolic ratio in dimension three” , Ann. Scuola Norm. Sup. Pisa Cl. Sci. (5) 19 (2019) no. 4, p. 1561-1582
\bibitem{Albach:2021wi} B. Albach \& H. Geiges - “Surfaces of section for Seifert fibrations” , Arnold Math. J. 7 (2021) no. 4, p. 573-597
\bibitem{Alvarez-Paiva:2014uq} J. C. Álvarez Paiva \& F. Balacheff - “Contact geometry and isosystolic inequalities” , Geom. Funct. Anal. 24 (2014) no. 2, p. 648-669
\bibitem{Banyaga:1978} A. Banyaga - “Sur la structure du groupe des difféomorphismes qui préservent une forme symplectique” , Comment. Math. Helv. 53 (1978) no. 2, p. 174-227
\bibitem{Banyaga:1997} A. Banyaga - The structure of classical diffeomorphism groups , Mathematics and its Appl. , vol. 400 , Kluwer Academic Publishers Group, Dordrecht, 1997
\bibitem{Besse:1978pr} A. L. Besse - Manifolds all of whose geodesics are closed , Ergeb. Math. Grenzgeb. , vol. 93 , Springer-Verlag, Berlin-New York, 1978
\bibitem{Benedetti:2021aa} G. Benedetti \& J. Kang - “A local contact systolic inequality in dimension three” , J. Eur. Math. Soc. (JEMS) 23 (2021) no. 3, p. 721-764
\bibitem{Bottkol:1980} M. Bottkol - “Bifurcation of periodic orbits on manifolds and Hamiltonian systems” , J. Differential Equations 37 (1980) no. 1, p. 12-22
\bibitem{Bialy:1994} M. Bialy \& L. Polterovich - “Geodesics of Hofer’s metric on the group of Hamiltonian diffeomorphisms” , Duke Math. J. 76 (1994) no. 1, p. 273-292
\bibitem{Calabi:1970} E. Calabi - “On the group of automorphisms of a symplectic manifold” , in Problems in analysis , Princeton University Press, Princeton, NJ, 1970, p. 1-26
\bibitem{Cristofaro-Gardiner:2020aa} D. Cristofaro-Gardiner \& M. Mazzucchelli - “The action spectrum characterizes closed contact 3-manifolds all of whose Reeb orbits are closed” , Comment. Math. Helv. 95 (2020) no. 3, p. 461-481
\bibitem{Croke:1994} C. B. Croke \& B. Kleiner - “Conjugacy and rigidity for manifolds with a parallel vector field” , J. Differential Geom. 39 (1994) no. 3, p. 659-680
\bibitem{Epstein:1972} D. B. A. Epstein - “Periodic flows on three-manifolds” , Ann. of Math. (2) 95 (1972), p. 66-82
\bibitem{Geiges:2020aa} H. Geiges - “What does a vector field know about volume?” , J. Fixed Point Theory Appl. 24 (2022) no. 2, article ID 23, 26 pages
\bibitem{Gromoll:1981kl} D. Gromoll \& K. Grove - “On metrics on $S^{2}$ all of whose geodesics are closed” , Invent. Math. 65 (1981/82) no. 1, p. 175-177
\bibitem{Ginzburg:1987lq} V. L. Ginzburg - “New generalizations of Poincaré’s geometric theorem” , Funktsional. Anal. i Prilozhen. 21 (1987) no. 2, p. 16-22
\bibitem{Giroux:2020} E. Giroux - “Ideal Liouville domains, a cool gadget” , J. Symplectic Geom. 18 (2020) no. 3, p. 769-790
\bibitem{Geiges:2018aa} H. Geiges \& C. Lange - “Seifert fibrations of lens spaces” , Abh. Math. Sem. Univ. Hamburg 88 (2018) no. 1, p. 1-22, Correction: Ibid. 91 (2021), no. 1, p. 145–150
\bibitem{Jankins:1983zm} M. Jankins \& W. D. Neumann - Lectures on Seifert manifolds , Brandeis Lect. Notes , vol. 2 , Brandeis University, Waltham, MA, 1983
\bibitem{Katok:1973mw} A. B. Katok - “Ergodic perturbations of degenerate integrable Hamiltonian systems” , Izv. Akad. Nauk SSSR Ser. Mat. 37 (1973), p. 539-576
\bibitem{Kegel:2021vo} M. Kegel \& C. Lange - “A Boothby-Wang theorem for Besse contact manifolds” , Arnold Math. J. 7 (2021) no. 2, p. 225-241
\bibitem{Lange:2020} C. Lange - “On metrics on 2-orbifolds all of whose geodesics are closed” , J. reine angew. Math. 758 (2020), p. 67-94
\bibitem{Lalonde:1995} F. Lalonde \& D. McDuff - “Hofer’s $L^\infty $-geometry: energy and stability of Hamiltonian flows. I, II” , Invent. Math. 122 (1995) no. 1, p. 1-33, 35–69, Erratum: Ibid. 123 (1996), no. 3, p. 613
\bibitem{Lisca:2004oz} P. Lisca \& G. Matić - “Transverse contact structures on Seifert 3-manifolds” , Algebraic Geom. Topol. 4 (2004), p. 1125-1144
\bibitem{Lange:2021} C. Lange \& T. Soethe - “Sharp systolic inequalities for rotationally symmetric $2$-orbifolds” , 2021
\bibitem{Mazzucchelli:2020aa} M. Mazzucchelli \& M. Radeschi - “On the structure of Besse convex contact spheres” , 2020
\bibitem{McDuff:1998uu} D. McDuff \& D. Salamon - Introduction to symplectic topology , Oxford Graduate Texts in Math. , Oxford University Press, Oxford, 2017
\bibitem{Orlik:1972ts} P. Orlik - Seifert manifolds , Lect. Notes in Math. , vol. 291 , Springer-Verlag, Berlin-New York, 1972
\bibitem{Radeschi:2017dz} M. Radeschi \& B. Wilking - “On the Berger conjecture for manifolds all of whose geodesics are closed” , Invent. Math. 210 (2017) no. 3, p. 911-962
\bibitem{Saglam:2018tb} M. Sağlam - “Contact forms with large systolic ratio in arbitrary dimensions” , Ann. Scuola Norm. Sup. Pisa Cl. Sci. (5) 22 (2021) no. 3, p. 1265-1308
\bibitem{Sullivan:1978} D. Sullivan - “A foliation of geodesics is characterized by having no ‘tangent homologies’” , J. Pure Appl. Algebra 13 (1978) no. 1, p. 101-104
\bibitem{Taubes:2007wi} C. H. Taubes - “The Seiberg-Witten equations and the Weinstein conjecture” , Geom. Topol. 11 (2007), p. 2117-2202
\bibitem{Wadsley:1975sp} A. W. Wadsley - “Geodesic foliations by circles” , J. Differential Geom. 10 (1975) no. 4, p. 541-549
\bibitem{Weinstein:1973} A. Weinstein - “Normal modes for nonlinear Hamiltonian systems” , Invent. Math. 20 (1973), p. 47-57
\bibitem{Ziller:1983rw} W. Ziller - “Geometry of the Katok examples” , Ergodic Theory Dynam. Systems 3 (1983) no. 1, p. 135-157
\end{thebibliography}
\end{document}
